% Oppsett for PDF-kompendiene fra fysikkhjelp.no (brukes av lag_pdf.py via Quarto/xelatex).
% Farger og bokser som i TRES0312-kompendiet; skrift og forside som på nettsiden.
\usepackage{fontspec}
\setmainfont{ComicNeue}[Path=FONTSTI/, Extension=.ttf, UprightFont=*-Regular, BoldFont=*-Bold, ItalicFont=*-Italic, BoldItalicFont=*-BoldItalic]
\newfontfamily\display{Montserrat}[Path=FONTSTI/, Extension=.ttf, UprightFont=*-ExtraBold, BoldFont=*-Black]
\newfontfamily\kritt{ComicNeue}[Path=FONTSTI/, Extension=.ttf, UprightFont=*-Bold]
\usepackage{xcolor,tikz,graphicx,amsmath,amssymb}
\usetikzlibrary{decorations.pathmorphing,arrows.meta,calc}
\usepackage[most]{tcolorbox}
\usepackage{qrcode}
\usepackage{newunicodechar}
\newunicodechar{✓}{\ensuremath{\checkmark}}
\newunicodechar{−}{\ensuremath{-}}
\newunicodechar{→}{\ensuremath{\rightarrow}}
\newunicodechar{≈}{\ensuremath{\approx}}

% Nettsidens farger (forsiden) og kompendiets farger (innmaten)
\definecolor{slate}{HTML}{0d1110}\definecolor{line}{HTML}{2a3330}
\definecolor{chalk}{HTML}{ecefe6}\definecolor{dust}{HTML}{9aa49e}\definecolor{neon}{HTML}{6ef04a}
\definecolor{cyan}{HTML}{56d7b8}\definecolor{orange}{HTML}{f19c2c}
\definecolor{headCol}{RGB}{18,82,88}
\definecolor{exmscol}{HTML}{e04b52}
\definecolor{oppgcol}{HTML}{1A6B9A}
\definecolor{defscol}{HTML}{449944}
\definecolor{graabg}{HTML}{F1F1EE}
\definecolor{losbg}{HTML}{FBF6F5}

% Overskrifter
\addtokomafont{disposition}{\color{headCol}}
\addtokomafont{chapter}{\kritt\Huge}
\addtokomafont{chapterprefix}{\display\normalsize\color{exmscol}}
\renewcommand*{\chapterformat}{\mbox{\display\large\color{exmscol}KAPITTEL \thechapter}}
\RedeclareSectionCommand[beforeskip=0pt, afterskip=1.6\baselineskip, innerskip=4mm]{chapter}
\RedeclareSectionCommand[beforeskip=-2.2\baselineskip, afterskip=1\baselineskip]{section}

% Topp- og bunntekst
\usepackage[headsepline=.4pt]{scrlayer-scrpage}
\clearpairofpagestyles
\ihead{\small\color{dust!60!black}fysikk\textcolor{defscol}{hjelp}.no \enspace·\enspace FAGTITTEL}
\ohead{\small\color{dust!60!black}\leftmark}
\ofoot*{\small\pagemark}
\automark[chapter]{chapter}

% Bokser
\tcbset{boks/.style={enhanced, breakable, boxrule=0pt, arc=2pt, left=10pt, right=8pt, top=6pt, bottom=6pt,
  before skip=10pt, after skip=10pt}}
\newtcolorbox{eksempel}[1]{boks, colback=white, borderline west={3pt}{0pt}{exmscol},
  title={\textit{#1}}, coltitle=exmscol, fonttitle=\bfseries, attach title to upper={\par\smallskip}}
\newtcolorbox{losningdel}{boks, colback=losbg, borderline west={2pt}{0pt}{exmscol!45}, before skip=6pt,
  title={Løsning}, coltitle=exmscol!80!black, fonttitle=\bfseries\small, attach title to upper={\par}}
\newtcolorbox{oppgave}[2]{boks, colback=oppgcol!5, borderline west={3pt}{0pt}{oppgcol},
  title={#1}, coltitle=oppgcol, fonttitle=\bfseries, attach title to upper={\hfill{\small\color{oppgcol!70}#2}\par\smallskip}}
\newtcolorbox{merkboks}[1]{boks, colback=headCol!4, borderline west={3pt}{0pt}{headCol!70}, title={#1}, coltitle=headCol, fonttitle=\bfseries, attach title to upper={\par\smallskip}}
\newtcolorbox{graaboks}{boks, colback=graabg, borderline west={3pt}{0pt}{defscol}}

% Video: QR-kode og lenke
\newcommand{\videoqr}[1]{\par\smallskip\noindent
  \begin{tcolorbox}[boks, colback=white, borderline={.4pt}{0pt}{line!40}, width=.78\linewidth]
    \begin{minipage}[c]{2.3cm}\qrcode[height=2cm]{#1}\end{minipage}\hfill
    \begin{minipage}[c]{\dimexpr\linewidth-2.7cm}{\bfseries\color{headCol}Video}\\[1mm]
      {\small Skann koden eller gå til}\\{\small\url{#1}}\end{minipage}
  \end{tcolorbox}\par}

\setlength{\parskip}{.45\baselineskip}
\setlength{\parindent}{0pt}
\usepackage{etoolbox}
\AtBeginEnvironment{oppgave}{\setlength{\parskip}{.3\baselineskip}}
